\documentclass[11pt,a4paper,UTF8]{ctexart}
\usepackage{geometry}
\geometry{left=18mm,right=18mm,top=24mm,bottom=24mm,headheight=22pt,footskip=12mm}
\usepackage{amsmath,amssymb,mathtools,bm,esint,extarrows,centernot}
\usepackage{xcolor}
\usepackage{tcolorbox}
\tcbuselibrary{skins,breakable}
\usepackage{fancyhdr}
\usepackage{titlesec}
\usepackage{hyperref}
\usepackage{tikz}
\usepackage{booktabs}
\usepackage{tabularx}

% --- Color Palette (Purple & DeepPink Academic Theme) ---
\definecolor{DeepPurple}{HTML}{4C1D95}
\definecolor{TitlePurple}{HTML}{581C87}
\definecolor{SubPurple}{HTML}{6B21A8}
\definecolor{DeepPink}{HTML}{FF1493}
\definecolor{LightPinkBg}{HTML}{FFF5F8}
\definecolor{LightPurpleBg}{HTML}{F8F5FF}
\definecolor{GoldAccent}{HTML}{D97706}
\definecolor{DarkText}{HTML}{1F2937}

\hypersetup{
    colorlinks=true,
    linkcolor=TitlePurple,
    citecolor=DeepPink,
    urlcolor=SubPurple,
    pdfborder={0 0 0}
}

% --- Typography & Spacing ---
\linespread{1.3}
\setlength{\parskip}{0.35em plus 0.1em minus 0.1em}
\setlength{\parindent}{0pt}

% --- Section Titles Styling ---
\titleformat{\section}
  {\Large\bfseries\color{TitlePurple}}{\thesection}{1em}{}
  [\vspace{1mm}{\color{TitlePurple!30}\hrule height 1pt}]
\titleformat{\subsection}
  {\large\bfseries\color{SubPurple}}{\thesubsection}{0.8em}{}
\titleformat{\subsubsection}
  {\normalsize\bfseries\color{DeepPurple}}{\thesubsubsection}{0.6em}{}

% --- tcolorbox Environments ---
% 1. Theorem / Formula / Method / Analysis Box (Deep Purple)
\newtcolorbox{thmbox}[1][]{
    enhanced,
    breakable,
    colback=LightPurpleBg,
    colframe=TitlePurple,
    arc=3pt,
    boxrule=0pt,
    leftrule=3.5pt,
    left=10pt,
    right=10pt,
    top=8pt,
    bottom=8pt,
    title=#1,
    coltitle=white,
    colbacktitle=TitlePurple,
    fonttitle=\bfseries\small,
    attach boxed title to top left={yshift=-2mm, xshift=4mm},
    boxed title style={boxrule=0pt, arc=2pt}
}

% 2. Solution / Formula / Derivation Box (DeepPink)
\newtcolorbox{solbox}[1][]{
    enhanced,
    breakable,
    colback=LightPinkBg,
    colframe=DeepPink,
    arc=3pt,
    boxrule=0pt,
    leftrule=3.5pt,
    left=10pt,
    right=10pt,
    top=8pt,
    bottom=8pt,
    title=#1,
    coltitle=white,
    colbacktitle=DeepPink,
    fonttitle=\bfseries\small,
    attach boxed title to top left={yshift=-2mm, xshift=4mm},
    boxed title style={boxrule=0pt, arc=2pt}
}

% 3. Learning Goals / Summary Box
\newtcolorbox{targetbox}[1][]{
    enhanced,
    breakable,
    colback=LightPurpleBg,
    colframe=DeepPurple,
    arc=3pt,
    boxrule=1pt,
    left=10pt,
    right=10pt,
    top=8pt,
    bottom=8pt,
    title=#1,
    coltitle=white,
    colbacktitle=DeepPurple,
    fonttitle=\bfseries\small,
    attach boxed title to top left={yshift=-2mm, xshift=4mm},
    boxed title style={boxrule=0pt, arc=2pt}
}

\begin{document}

% ------------------- 讲义封面 (Cover Page) -------------------
\begin{titlepage}
\thispagestyle{empty}
\begin{tikzpicture}[remember picture,overlay]
    \fill[DeepPurple] (current page.north west) rectangle ([yshift=-7cm]current page.north east);
    \draw[DeepPink, line width=3.5pt] ([yshift=-7cm]current page.north west) -- ([yshift=-7cm]current page.north east);
    \draw[GoldAccent, line width=1pt] ([yshift=-7.12cm]current page.north west) -- ([yshift=-7.12cm]current page.north east);
    
    \node[anchor=north east, opacity=0.12, text=white] at ([xshift=-1.5cm, yshift=-0.5cm]current page.north east) {
        \fontsize{70}{70}\selectfont\bfseries 2027
    };
\end{tikzpicture}

\vspace*{0.8cm}
\begin{center}
    {\color{white}\fontsize{26}{32}\selectfont\bfseries 2027 考研数学(二) 核心讲义}\\[0.6cm]
    {\color{white}\fontsize{13}{18}\selectfont\bfseries 全国硕士研究生招生考试 · 统考数学(二) 核心精讲全书}\\[1.5cm]
\end{center}

\vspace{3.5cm}
\begin{center}
    {\color{GoldAccent}\large\bfseries $\blacktriangleright$ 基础强化 · 核心题解 · 考点深水区 $\blacktriangleleft$}\\[0.8cm]
    {\color{TitlePurple}\fontsize{24}{28}\selectfont\bfseries 第六章 一元函数微分学的应用(中值定理\&微分等式\&微分不等式)}\\[0.6cm]
    {\color{DeepPink}\large\bfseries —— 核心精讲大例题与题解三步法 ——}\\[1.5cm]
\end{center}

\vfill

\begin{center}
\begin{tcolorbox}[
    colback=white,
    colframe=DeepPurple,
    width=0.88\textwidth,
    arc=4pt,
    boxrule=1.2pt,
    halign=center,
    drop shadow=gray!30
]
    \vspace{3mm}
    {\bfseries\large 编著团队：EscoffierZhou}\\[2.5mm]
    {\bfseries\color{TitlePurple} 目标院校：中国科学院大学 (UCAS) 计算机专硕 (22408)}\\[2.5mm]
    {\small\color{gray} 备考铁律：手算真题数字 · 错题白纸重做 · 408 客观题为王}\\[2.5mm]
    {\small 编制版本：2027 届首发版 · \today}
    \vspace{2mm}
\end{tcolorbox}
\end{center}
\vspace{0.8cm}
\end{titlepage}

% ------------------- 目录与页面主体配置 -------------------
\pagestyle{fancy}
\fancyhf{}
\fancyhead[L]{\small\color{gray}2027 考研数学(二) · 核心讲义系列}
\fancyhead[R]{\small\color{TitlePurple}\nouppercase{\leftmark}}
\fancyfoot[C]{\small\color{gray}— 第 \thepage\ 页 —}
\fancyfoot[R]{\small\color{gray}UCAS 22408}
\renewcommand{\headrulewidth}{0.5pt}

\pagenumbering{Roman}
\tableofcontents
\newpage
\pagenumbering{arabic}


\section{6.例题}


\subsection{例题 1: 零点定理辅助函数与端点变号}

设函数$f(x)$在$[0, 1]$上连续，且$f(1) = 0$，$f\left(\frac{1}{2}\right) = 1$。证明：存在$\eta \in \left(\frac{1}{2}, 1\right)$，使得$f(\eta) = \eta$。

\textbf{\color{DeepPurple}【思路点拨】} 结论等价于$f(\eta) - \eta = 0$，构造辅助函数$F(x) = f(x) - x$，利用在区间端点处的正负变号由零点定理证得结论


\begin{solbox}[分步推导与答题规范]
\textbf{[1]} 构造辅助函数：

待证结论等价于$f(\eta) - \eta = 0$，故构造辅助函数：


\begin{equation*}
F(x) = f(x) - x
\end{equation*}


\textbf{[2]} 连续性与端点值计算：

因为$f(x)$在$[0, 1]$上连续，所以$F(x)$在闭区间$\left[\frac{1}{2}, 1\right]$上连续。

计算端点函数值：

$F\left(\frac{1}{2}\right) = f\left(\frac{1}{2}\right) - \frac{1}{2} = 1 - \frac{1}{2} = \frac{1}{2} > 0$

$F(1) = f(1) - 1 = 0 - 1 = -1 < 0$。

\textbf{[3]} 应用零点定理：

由于$F\left(\frac{1}{2}\right) \cdot F(1) < 0$，由闭区间连续函数零点定理可知：

在开区间$\left(\frac{1}{2}, 1\right)$内至少存在一点$\eta$，使得$F(\eta) = 0$，即$f(\eta) = \eta$。
\end{solbox}


\subsection{例题 2: 极限局部保号性与零点定理}

设函数$f(x)$在区间$[0, 1]$上连续，且$f(1) > 0$，$\lim_{x \to 0^+} \frac{f(x)}{x} < 0$。证明：方程$f(x) = 0$在区间$(0, 1)$内至少存在一个实根。

\textbf{\color{DeepPurple}【思路点拨】} 利用极限局部保号性在原点右邻域内锁定一个函数值为负的点，再与右端点$f(1) > 0$配合应用零点定理


\begin{solbox}[分步推导与答题规范]
\textbf{[1]} 极限保号性分析：

记极限值$A = \lim_{x \to 0^+} \frac{f(x)}{x} < 0$。由极限局部保号性，存在$\delta \in (0, 1)$，当$x \in (0, \delta)$时，恒有：


\begin{equation*}
\frac{f(x)}{x} < 0
\end{equation*}


\textbf{[2]} 取点锁定负值：

任取点$a \in (0, \delta)$，由于$a > 0$，必有$f(a) < 0$。

\textbf{[3]} 应用零点定理：

已知$f(1) > 0$，且$f(x)$在闭区间$[a, 1]$上连续，满足：


\begin{equation*}
f(a) \cdot f(1) < 0
\end{equation*}


根据零点定理，在开区间$(a, 1) \subset (0, 1)$内至少存在一点$\xi$，使得$f(\xi) = 0$，即方程在$(0, 1)$内至少有一个实根。
\end{solbox}


\subsection{例题 3: 导数零点定理(达布定理)证明}

设$f(x)$在$[a, b]$上可导，证明：当$f'_+(a) \cdot f'_-(b) < 0$时，存在$\xi \in (a, b)$，使得$f'(\xi) = 0$。

\textbf{\color{DeepPurple}【思路点拨】} 利用导数定义和保号性排除端点为最值点的可能性（脱帽法），结合闭区间连续函数最值定理说明最值必在开区间内部取得，最后由费马引理证得导数为零


\begin{solbox}[分步推导与答题规范]
\textbf{[1]} 排除端点最值（不妨设$f'_+(a) > 0$，$f'_-(b) < 0$）：

由右导数定义及保号性：


\begin{equation*}
f'_+(a) = \lim_{x \to a^+} \frac{f(x) - f(a)}{x - a} > 0 \implies \exists \delta_1 > 0, \text{ 当 } x \in (a, a+\delta_1) \text{ 时, } f(x) > f(a)
\end{equation*}


说明左端点$f(a)$绝不是$[a, b]$上的最大值。

同理由左导数定义及保号性：


\begin{equation*}
f'_-(b) = \lim_{x \to b^-} \frac{f(x) - f(b)}{x - b} < 0 \implies \exists \delta_2 > 0, \text{ 当 } x \in (b-\delta_2, b) \text{ 时, } f(x) > f(b)
\end{equation*}


说明右端点$f(b)$也绝不是$[a, b]$上的最大值。

\textbf{[2]} 最值必在区间内部取得：

因为$f(x)$在$[a, b]$上可导必连续，由闭区间连续函数最值定理，$f(x)$在$[a, b]$上必能取得最大值。

由于两个端点均取不到最大值，该最大值必在开区间内部某点$\xi \in (a, b)$处取得。

\textbf{[3]} 费马引理得出结论：

因为$\xi \in (a, b)$是内部极值点且$f(x)$在$\xi$处可导，由费马引理必有：


\begin{equation*}
f'(\xi) = 0
\end{equation*}

\end{solbox}


\subsection{例题 4: 离散平均值定理与罗尔定理串联}

设$f(x)$在$[0, 3]$上连续，在$(0, 3)$内可导，且$f(0) + f(1) + f(2) = 3$，$f(3) = 1$。证明：必存在$\xi \in (0, 3)$，使得$f'(\xi) = 0$。

\textbf{\color{DeepPurple}【思路点拨】} 利用连续函数最值定理与介值定理先找到一个函数值等于平均值$1$的点$c$，再在$[c, 3]$上应用罗尔定理


\begin{solbox}[分步推导与答题规范]
\textbf{[1]} 由介值定理寻找平均值点：

$f(x)$在闭区间$[0, 2]$上连续，由最值定理必有最小值$m$和最大值$M$，故：


\begin{equation*}
m \le f(0) \le M, \quad m \le f(1) \le M, \quad m \le f(2) \le M
\end{equation*}


三式相加除以3：


\begin{equation*}
m \le \frac{f(0) + f(1) + f(2)}{3} \le M \implies m \le 1 \le M
\end{equation*}


由介值定理，至少存在一点$c \in [0, 2]$，使得：


\begin{equation*}
f(c) = \frac{f(0) + f(1) + f(2)}{3} = 1
\end{equation*}


\textbf{[2]} 应用罗尔定理：

考察区间$[c, 3]$：由于$c \le 2 < 3$，区间非空。满足：

- $f(x)$在$[c, 3]$上连续，在$(c, 3)$内可导；
- 端点值相等：$f(c) = 1 = f(3)$。

由罗尔定理，必存在$\xi \in (c, 3) \subset (0, 3)$，使得$f'(\xi) = 0$。
\end{solbox}


\subsection{例题 5: 微分方程还原法构造罗尔辅助函数}

设$f(x)$在$[a, b]$上连续，在$(a, b)$内可导，$f(a) = 0$，$a > 0$。证明：存在$\xi \in (a, b)$，使得$f(\xi) = \frac{b - \xi}{a} f'(\xi)$。

\textbf{\color{DeepPurple}【思路点拨】} 将待证等式整理为微分方程形式$(b - x)f'(x) - a f(x) = 0$，通过积分因子法求得辅助函数$F(x) = f(x)(b - x)^a$，验证端点后应用罗尔定理


操作[1]:微分方程分离变量与求原函数：

\begin{solbox}[分步推导与答题规范]
将欲证等式整理为：


\begin{equation*}
(b - \xi)f'(\xi) - a f(\xi) = 0 \iff \frac{f'(\xi)}{f(\xi)} - \frac{a}{b - \xi} = 0
\end{equation*}


两边积分寻找原函数：


\begin{equation*}
\int \left(\frac{f'}{f} - \frac{a}{b - x}\right)\mathrm{d}x = \ln \vert{}f(x)\vert{} + a \ln (b - x) = \ln \left[\vert{}f(x)\vert{} (b - x)^a\right]
\end{equation*}


由此锁定辅助函数为：
\begin{equation*}
F(x) = f(x)(b - x)^a
\end{equation*}

\end{solbox}

操作[2]:验证罗尔定理条件：

\begin{solbox}[分步推导与答题规范]
- $F(a) = f(a)(b - a)^a = 0 \cdot (b - a)^a = 0$；
- $F(b) = f(b)(b - b)^a = 0$。

$F(x)$在$[a, b]$上连续，在$(a, b)$内可导，且$F(a) = F(b) = 0$。
\end{solbox}

操作[3]:应用罗尔定理：

\begin{solbox}[分步推导与答题规范]
由罗尔定理，必存在$\xi \in (a, b)$，使得$F'(\xi) = 0$。

求导展开：


\begin{equation*}
F'(x) = f'(x)(b - x)^a - a f(x)(b - x)^{a-1} = (b - x)^{a-1} \left[(b - x)f'(x) - a f(x)\right]
\end{equation*}


因为$\xi \in (a, b) \implies b - \xi \neq 0$，故由$F'(\xi) = 0$可得：


\begin{equation*}
(b - \xi)f'(\xi) - a f(\xi) = 0 \implies f(\xi) = \frac{b - \xi}{a} f'(\xi)
\end{equation*}

\end{solbox}


\subsection{例题 6: 三点共线与二次罗尔定理}

设函数$f(x)$在$[0, 1]$上连续，在$(0, 1)$内二阶可导，过点$A(0, f(0))$与点$B(1, f(1))$的直线与曲线$y = f(x)$相交于点$C(c, f(c))$，其中$0 < c < 1$。证明：存在$\xi \in (0, 1)$，使得$f''(\xi) = 0$。

\textbf{\color{DeepPurple}【思路点拨】} 构造曲线与割线之差的辅助函数$F(x) = f(x) - y_1(x)$，三点共线提供三个零点，连续两次应用罗尔定理证明二阶导为零


操作[1]:设割线并构造差函数：

\begin{solbox}[分步推导与答题规范]
连接$A, B$两点的割线方程为：


\begin{equation*}
y_1(x) = f(0) + \frac{f(1) - f(0)}{1 - 0}(x - 0) = [f(1) - f(0)]x + f(0)
\end{equation*}


构造辅助函数（曲线与割线之差）：


\begin{equation*}
F(x) = f(x) - y_1(x) = f(x) - [f(1) - f(0)]x - f(0)
\end{equation*}

\end{solbox}

操作[2]:验证特征零点：

\begin{solbox}[分步推导与答题规范]
- $F(0) = f(0) - f(0) = 0$；
- $F(1) = f(1) - [f(1) - f(0)] - f(0) = 0$；
- 因为点$C(c, f(c))$在割线上，割线纵坐标$y_1(c) = f(c)$，故$F(c) = 0$。
\end{solbox}

操作[3]:两次应用罗尔定理：

\begin{solbox}[分步推导与答题规范]
- 在$[0, c]$上由$F(0) = F(c) = 0$，存在$\xi_1 \in (0, c)$使得$F'(\xi_1) = 0$；
- 在$[c, 1]$上由$F(c) = F(1) = 0$，存在$\xi_2 \in (c, 1)$使得$F'(\xi_2) = 0$；
- 显然$\xi_1 < \xi_2$。在$[\xi_1, \xi_2]$上对$F'(x)$应用罗尔定理，存在$\xi \in (\xi_1, \xi_2) \subset (0, 1)$，使得$F''(\xi) = 0$。
\end{solbox}

操作[4]:还原原函数：

\begin{solbox}[分步推导与答题规范]
因为$y_1(x)$是一次多项式，$y_1''(x) \equiv 0$，故$F''(x) = f''(x)$。

从而代入即得$f''(\xi) = 0$。
\end{solbox}


\subsection{例题 7: 复合导数结构的两层罗尔定理链}

设函数$f(x)$在区间$[0, 1]$上具有二阶导数，且$f(1) > 0$，$\lim_{x \to 0^+} \frac{f(x)}{x} < 0$。证明：方程$f(x)f''(x) + [f'(x)]^2 = 0$在区间$(0, 1)$内至少存在两个不同实根。

\textbf{\color{DeepPurple}【思路点拨】} 识别待证式为$[f(x)f'(x)]' = 0$，构造辅助函数$F(x) = f(x)f'(x)$；利用极限与零点定理找出$F(x)$的3个零点，两次应用罗尔定理证得两个不同实根


操作[1]:识别辅助函数结构：

\begin{solbox}[分步推导与答题规范]
观察待证方程，左端可直接还原为积的求导法则：


\begin{equation*}
[f(x)f'(x)]' = f(x)f''(x) + [f'(x)]^2
\end{equation*}


由此设辅助函数为：
\begin{equation*}
F(x) = f(x)f'(x)
\end{equation*}

\end{solbox}

操作[2]:寻找零点：

\begin{solbox}[分步推导与答题规范]
- 由极限$\lim_{x \to 0^+} \frac{f(x)}{x}$存在且分母趋于0，由分子趋于0知$f(0) = 0$；
- 由例题[2]的结论，因$\lim_{x \to 0^+} \frac{f(x)}{x} < 0$且$f(1) > 0$，必存在$b \in (0, 1)$使得$f(b) = 0$；
- 在区间$[0, b]$上，$f(0) = f(b) = 0$，由罗尔定理，存在$c \in (0, b)$使得$f'(c) = 0$。
\end{solbox}

操作[3]:验证$F(x)$的三个零点：

\begin{solbox}[分步推导与答题规范]
- $F(0) = f(0)f'(0) = 0$；
- $F(c) = f(c)f'(c) = f(c) \cdot 0 = 0$；
- $F(b) = f(b)f'(b) = 0 \cdot f'(b) = 0$。

其中$0 < c < b < 1$。
\end{solbox}

操作[4]:分段应用罗尔定理：(证明 $\xi \neq \eta$)

\begin{solbox}[分步推导与答题规范]
- 在$[0, c]$上对$F(x)$应用罗尔定理，存在$\xi \in (0, c)$使得$F'(\xi) = 0$；
- 在$[c, b]$上对$F(x)$应用罗尔定理，存在$\eta \in (c, b)$使得$F'(\eta) = 0$。

显然$0 < \xi < c < \eta < b < 1$，即$\xi \neq \eta$。

故原方程在$(0, 1)$内至少存在两个不同实根。
\end{solbox}


\subsection{例题 8: 导数有界判定函数有界(拉格朗日中值定理)}

若$f(x)$在$(a, b)$内可导且$f'(x)$有界，证明：$f(x)$在$(a, b)$内有界。

\textbf{\color{DeepPurple}【思路点拨】} 固定区间内一个参考点$x_0$，任取$x$应用拉格朗日中值定理将增量转化为导数控制，结合绝对值三角不等式放缩完成证明


\begin{solbox}[分步推导与答题规范]
\textbf{[1]} 固定参考点并应用拉格朗日中值定理：

在开区间$(a, b)$内任取一固定参考点$x_0 \in (a, b)$，则$f(x_0)$为确定常数。

对任意$x \in (a, b)$，在以$x$和$x_0$为端点的闭区间上应用拉格朗日中值定理：


\begin{equation*}
f(x) - f(x_0) = f'(\xi)(x - x_0) \quad (\xi \text{ 介于 } x \text{ 与 } x_0 \text{ 之间})
\end{equation*}


\textbf{[2]} 绝对值放缩：

移项并取绝对值：


\begin{equation*}
\vert{}f(x)\vert{} = \vert{}f(x_0) + f'(\xi)(x - x_0)\vert{} \le \vert{}f(x_0)\vert{} + \vert{}f'(\xi)\vert{} \cdot \vert{}x - x_0\vert{}
\end{equation*}


\textbf{[3]} 有界性判定：

已知$f'(x)$在$(a, b)$上有界，即存在常数$k > 0$，使得对任意$\xi \in (a, b)$均有$\vert{}f'(\xi)\vert{} \le k$；

又显然$\vert{}x - x_0\vert{} < b - a$。

从而有：


\begin{equation*}
\vert{}f(x)\vert{} \le \vert{}f(x_0)\vert{} + k(b - a) \triangleq M
\end{equation*}


由于$M$为与$x$无关的常数，因此$f(x)$在$(a, b)$内有界。
\end{solbox}


\subsection{例题 9: 函数平方求导构造与拉格朗日中值定理}

设$f(x)$在$[0, 1]$上连续，在$(0, 1)$内可导，$f(0) = 0$，且$f(x)$在$[0, 1]$上不恒等于零。证明：存在$\xi \in (0, 1)$，使得$f(\xi)f'(\xi) > 0$。

\textbf{\color{DeepPurple}【思路点拨】} 欲证$f(\xi)f'(\xi) > 0$，联想$F(x) = \frac{1}{2}f^2(x)$的求导结构，由不恒为零锁定正值点，应用拉格朗日中值定理


\begin{solbox}[分步推导与答题规范]
\textbf{[1]} 构造辅助函数：

观察待证式$f(\xi)f'(\xi) > 0$，联想到平方函数的复合求导结构：


\begin{equation*}
F(x) = \frac{1}{2}f^2(x) \implies F'(x) = f(x)f'(x)
\end{equation*}


计算原点值：$F(0) = \frac{1}{2}f^2(0) = 0$。

\textbf{[2]} 锁定正值点：

因为$f(x)$在$[0, 1]$上不恒等于零，必存在某个点$a \in (0, 1]$，使得$f(a) \neq 0$。

从而必有：


\begin{equation*}
F(a) = \frac{1}{2}f^2(a) > 0
\end{equation*}


\textbf{[3]} 应用拉格朗日中值定理：

在闭区间$[0, a]$上对$F(x)$应用拉格朗日中值定理，存在$\xi \in (0, a) \subset (0, 1)$，使得：


\begin{equation*}
F'(\xi) = \frac{F(a) - F(0)}{a - 0} = \frac{\frac{1}{2}f^2(a) - 0}{a} > 0
\end{equation*}


将$F'(\xi) = f(\xi)f'(\xi)$代入，即证存在$\xi \in (0, 1)$，满足$f(\xi)f'(\xi) > 0$。
\end{solbox}


\subsection{例题 10: 割线斜率构造与商函数单调性}

设函数$f(x)$满足$f(0) = 0$，且当$x > 0$时，$f(x) < 0$，$f'(x) < 0$，$f''(x) > 0$。则当$0 < a < x < b$时，有( )

(A) $xf(x) > af(a)$ $\qquad$(B) $bf(b) > xf(x)$

(C) $xf(a) > af(x)$ $\qquad$(D) $xf(b) > bf(x)$

\textbf{\color{DeepPurple}【思路点拨】} 两边同除转化为判断商函数$h(x) = \frac{f(x)}{x}$的单调性，对分子用拉格朗日中值定理及二阶导正号证明单调递增


\begin{solbox}[分步推导与答题规范]
\textbf{[1]} 转化为商函数单调性：

欲比较形如$xf(a)$与$af(x)$的大小，两边同除以正数$ax$，等价于考察函数$h(x) = \frac{f(x)}{x}\ (x > 0)$的单调性。

\textbf{[2]} 求导与分子化简：


\begin{equation*}
h'(x) = \frac{xf'(x) - f(x)}{x^2}
\end{equation*}


对分子中的$f(x)$应用拉格朗日中值定理（利用$f(0) = 0$）：


\begin{equation*}
f(x) = f(x) - f(0) = f'(\xi)(x - 0) = xf'(\xi) \quad (0 < \xi < x)
\end{equation*}


代入导数分子中：


\begin{equation*}
xf'(x) - f(x) = xf'(x) - xf'(\xi) = x[f'(x) - f'(\xi)]
\end{equation*}


\textbf{[3]} 利用二阶导数判定符号：

因为$f''(x) > 0$，所以导函数$f'(x)$在$(0, +\infty)$上严格单调增加。

由$0 < \xi < x$必有$f'(\xi) < f'(x) \implies f'(x) - f'(\xi) > 0$。

又$x > 0$，故分子严格大于0，从而$h'(x) > 0$恒成立，说明$h(x) = \frac{f(x)}{x}$在$(0, +\infty)$上严格单调增加。

\textbf{[4]} 区间不等式判定：

当$0 < a < x < b$时，必有：


\begin{equation*}
\frac{f(a)}{a} < \frac{f(x)}{x} < \frac{f(b)}{b}
\end{equation*}


- 由$\frac{f(a)}{a} < \frac{f(x)}{x} \implies xf(a) < af(x)$（排除 C）；
- 由$\frac{f(x)}{x} < \frac{f(b)}{b} \implies bf(x) < xf(b) \iff xf(b) > bf(x)$（符合 D）。

\textbf{[5]} 故正确选项为 \textbf{(D)}。
\end{solbox}


\subsection{例题 11: 柯西中值定理结构识别}

设$f(x)$在$[a, b]$上连续，在$(a, b)$内可导，$0 < a < b$。证明：至少存在一点$\xi \in (a, b)$，使得：


\begin{equation*}
f(b) - f(a) = \xi \ln\frac{b}{a} f'(\xi)
\end{equation*}

\textbf{\color{DeepPurple}【思路点拨】} 改写为比值形式$\frac{f(b)-f(a)}{\ln b - \ln a} = \frac{f'(\xi)}{\frac{1}{\xi}}$，构造辅助函数$g(x) = \ln x$应用柯西中值定理


\begin{solbox}[分步推导与答题规范]
\textbf{[1]} 改写为分式比值形式：

将待证结论两端整理变形为：


\begin{equation*}
\frac{f(b) - f(a)}{\ln b - \ln a} = \xi f'(\xi) = \frac{f'(\xi)}{\frac{1}{\xi}}
\end{equation*}


\textbf{[2]} 构造辅助函数与验证条件：

构造分母对应的函数$g(x) = \ln x$：

- $f(x)$与$g(x)$在$[a, b]$上均连续，在$(a, b)$内均可导；
- 当$x \in (a, b)$时，由$0 < a < b$知$g'(x) = \frac{1}{x} \neq 0$。

\textbf{[3]} 应用柯西中值定理：

至少存在一点$\xi \in (a, b)$，使得：


\begin{equation*}
\frac{f(b) - f(a)}{g(b) - g(a)} = \frac{f'(\xi)}{g'(\xi)}
\end{equation*}


将$g(b) - g(a) = \ln b - \ln a = \ln\frac{b}{a}$，以及$g'(\xi) = \frac{1}{\xi}$代入：


\begin{equation*}
\frac{f(b) - f(a)}{\ln\frac{b}{a}} = \frac{f'(\xi)}{\frac{1}{\xi}} = \xi f'(\xi)
\end{equation*}


两边同乘$\ln\frac{b}{a}$，即证：


\begin{equation*}
f(b) - f(a) = \xi \ln\frac{b}{a} f'(\xi)
\end{equation*}

\end{solbox}


\subsection{例题 12: 积分条件下的中点泰勒展开与积分保号}

设函数$f(x)$在$[0, 1]$上二阶可导，且$\int_0^1 f(x)\,\mathrm{d}x = 0$。则( )

(A) 当$f'(x) < 0$时，$f\left(\frac{1}{2}\right) < 0$ $\qquad$(B) 当$f''(x) < 0$时，$f\left(\frac{1}{2}\right) < 0$

(C) 当$f'(x) > 0$时，$f\left(\frac{1}{2}\right) < 0$ $\qquad$(D) 当$f''(x) > 0$时，$f\left(\frac{1}{2}\right) < 0$

\textbf{\color{DeepPurple}【思路点拨】} 在区间中点$x_0 = \frac{1}{2}$处作带拉格朗日余项的二阶泰勒展开，两端在$[0, 1]$积分，利用对称性消去一次项，由二阶导保号性定值


\textit{}看到 $\int_a^b f(x)\,\mathrm{d}x = 0$\textit{}

\begin{solbox}[分步推导与答题规范]
\textbf{第一反射（罗尔原函数）}：令变上限积分 $F(x) = \int_a^x f(t)\,\mathrm{d}t$

则 $F(a) = 0, F(b) = 0 \implies \exists \xi \in (a, b), F'(\xi) = f(\xi) = 0$（由积分等于 0 必能砸出一个原函数零点！）

\textbf{第二反射（积分中值定理）}：$\int_a^b f(x)\,\mathrm{d}x = f(\eta)(b - a) = 0 \implies f(\eta) = 0$。
\end{solbox}

\begin{solbox}[分步推导与答题规范]
\textbf{[1]} 在中点处进行泰勒展开：

将$f(x)$在$x_0 = \frac{1}{2}$处展开为带拉格朗日余项的泰勒公式：


\begin{equation*}
f(x) = f\left(\frac{1}{2}\right) + f'\left(\frac{1}{2}\right)\left(x - \frac{1}{2}\right) + \frac{f''(\xi)}{2}\left(x - \frac{1}{2}\right)^2 \quad (\xi \text{ 介于 } x \text{ 与 } \frac{1}{2} \text{ 之间})
\end{equation*}


\textbf{[2]} 两端在$[0, 1]$上求定积分：


\begin{equation*}
\int_0^1 f(x)\,\mathrm{d}x = \int_0^1 f\left(\frac{1}{2}\right)\mathrm{d}x + f'\left(\frac{1}{2}\right)\int_0^1 \left(x - \frac{1}{2}\right)\mathrm{d}x + \int_0^1 \frac{f''(\xi)}{2}\left(x - \frac{1}{2}\right)^2 \mathrm{d}x
\end{equation*}


\textbf{[3]} 利用对称性计算积分：

$\int_0^1 f\left(\frac{1}{2}\right)\mathrm{d}x = f\left(\frac{1}{2}\right)$；

一次奇对称项为零：$\int_0^1 \left(x - \frac{1}{2}\right)\mathrm{d}x = \left.\frac{1}{2}\left(x - \frac{1}{2}\right)^2\right\vert{}_0^1 = \frac{1}{8} - \frac{1}{8} = 0$。

代入$\int_0^1 f(x)\,\mathrm{d}x = 0$，移项整理得：


\begin{equation*}
f\left(\frac{1}{2}\right) = -\frac{1}{2}\int_0^1 f''(\xi)\left(x - \frac{1}{2}\right)^2 \mathrm{d}x
\end{equation*}


\textbf{[4]} 保号性判定：

当$f''(x) > 0$恒成立时，被积函数中$f''(\xi) > 0$且$\left(x - \frac{1}{2}\right)^2 \ge 0$（仅在一点为0），由定积分保号性可知积分值大于0。

因此必有$f\left(\frac{1}{2}\right) < 0$。

\textbf{[5]} 故正确选项为 \textbf{(D)}。
\end{solbox}


\subsection{例题 13: 高阶导数与对称点双向泰勒展开}

设函数$f(x)$在闭区间$[-1, 1]$上具有三阶连续导数，且$f(-1) = 0$，$f(1) = 1$，$f'(0) = 0$。证明：在开区间$(-1, 1)$内至少存在一点$\xi$，使得$f'''(\xi) = 3$。

\textbf{\color{DeepPurple}【思路点拨】} 在一阶导数已知点$x=0$处作带拉格朗日余项的二阶泰勒展开，分别代入$x = -1$与$x = 1$，两式相减消去未知的$f(0)$与$f''(0)$，结合介值定理证得结论


操作[1]:在导数已知点$x = 0$处作泰勒展开：

\begin{solbox}[分步推导与答题规范]
将$f(x)$在$x_0 = 0$处展开为带拉格朗日余项的泰勒公式：


\begin{equation*}
f(x) = f(0) + f'(0)x + \frac{1}{2}f''(0)x^2 + \frac{1}{6}f'''(\eta)x^3
\end{equation*}


代入$f'(0) = 0$化简得：


\begin{equation*}
f(x) = f(0) + \frac{1}{2}f''(0)x^2 + \frac{1}{6}f'''(\eta)x^3 \quad (\eta \text{ 介于 } 0 \text{ 与 } x \text{ 之间})
\end{equation*}

\end{solbox}

操作[2]:分别代入$x = -1$与$x = 1$：

\begin{solbox}[分步推导与答题规范]
- 取$x = -1$，存在$\eta_1 \in (-1, 0)$，使得：

 
\begin{equation*}
0 = f(-1) = f(0) + \frac{1}{2}f''(0) - \frac{1}{6}f'''(\eta_1) \tag{1}
\end{equation*}


- 取$x = 1$，存在$\eta_2 \in (0, 1)$，使得：

 
\begin{equation*}
1 = f(1) = f(0) + \frac{1}{2}f''(0) + \frac{1}{6}f'''(\eta_2) \tag{2}
\end{equation*}

\end{solbox}

操作[3]:相减消去低阶项：

\begin{solbox}[分步推导与答题规范]
式(2)减去式(1)：


\begin{equation*}
1 - 0 = \frac{1}{6}\left[f'''(\eta_1) + f'''(\eta_2)\right] \implies \frac{f'''(\eta_1) + f'''(\eta_2)}{2} = 3
\end{equation*}

\end{solbox}

操作[4]:应用连续函数介值定理：

\begin{solbox}[分步推导与答题规范]
因为$f(x)$具有三阶连续导数，$f'''(x)$在闭区间$[\eta_1, \eta_2] \subset (-1, 1)$上连续。

设$f'''(x)$在该区间上的最大值和最小值分别为$M$和$m$，则必有$m \le 3 \le M$。

由介值定理，至少存在一点$\xi \in [\eta_1, \eta_2] \subset (-1, 1)$，使得：


\begin{equation*}
f'''(\xi) = \frac{f'''(\eta_1) + f'''(\eta_2)}{2} = 3
\end{equation*}

\end{solbox}


\subsection{例题 14: 偶函数对称性降维与单调性求实根个数}

在区间$(-\infty, +\infty)$内，方程$\vert{}x\vert{}^{\frac{1}{2}} + \vert{}x\vert{}^{\frac{1}{3}} - \cos x = 0$( )

(A) 无实根 $\qquad$(B) 有且仅有一个实根 $\qquad$(C) 有且仅有两个实根 $\qquad$(D) 有无穷多个实根

\textbf{\color{DeepPurple}【思路点拨】} 利用偶函数对称性，只需分析正半轴$x > 0$；通过端点值与导数正号证明在$(0, 1)$上有且仅有一根，对称到全域得两实根


\begin{solbox}[分步推导与答题规范]
\textbf{[1]} 偶函数性质判定：

记$f(x) = \vert{}x\vert{}^{\frac{1}{2}} + \vert{}x\vert{}^{\frac{1}{3}} - \cos x$。

因为$f(-x) = f(x)$，故$f(x)$为\textbf{偶函数}，图像关于$y$轴对称。

代入$x = 0$：$f(0) = 0 + 0 - 1 = -1 \neq 0$，故$x = 0$不是根，只需研究$x > 0$上的根。

\textbf{[2]} 在$x > 0$侧缩小根的区间：

- 当$x \ge 1$时，$\vert{}x\vert{}^{\frac{1}{2}} \ge 1$且$\vert{}x\vert{}^{\frac{1}{3}} \ge 1$，而$\cos x \le 1$；

 故$f(x) \ge 1 + 1 - 1 = 1 > 0$恒成立，在$[1, +\infty)$内无根。

- 因此$f(x) = 0$的正实根只能落在区间$(0, 1)$内。

\textbf{[3]} 求导分析单调性：

当$x \in (0, 1)$时：


\begin{equation*}
f'(x) = \frac{1}{2}x^{-\frac{1}{2}} + \frac{1}{3}x^{-\frac{2}{3}} + \sin x
\end{equation*}


因各项均严格大于0，故$f'(x) > 0$恒成立，$f(x)$在$(0, 1)$上严格单调增加。

结合$f(0) = -1 < 0$，$f(1) \ge 1 > 0$，由零点定理与严格单调性可知，方程在$(0, 1)$内\textbf{有且仅有一个正实根}。

\textbf{[4]} 全域对称判定：

由偶函数对称性，在负半轴$(-\infty, 0)$上对应存在且仅存在一个对称负实根。

因此方程在全实数域内\textbf{有且仅有两个实根}。

\textbf{[5]} 故正确选项为 \textbf{(C)}。
\end{solbox}


\subsection{例题 15: 奇次多项式保底与二次判别法求唯一实根}

若$3a^2 - 5b < 0$，则方程$x^5 + 2ax^3 + 3bx + 4c = 0$( )

(A) 无实根 $\qquad$(B) 有唯一实根 $\qquad$(C) 有三个不同实根 $\qquad$(D) 有五个不同实根

\textbf{\color{DeepPurple}【思路点拨】} 奇次多项式两端极限异号至少有1根；对导函数换元为二次方程，由判别式小于零说明一阶导恒正，函数严格单增锁定唯一实根


\begin{solbox}[分步推导与答题规范]
\textbf{[1]} 实根存在性（保底1根）：

记$f(x) = x^5 + 2ax^3 + 3bx + 4c$。

因为最高次为5（奇次多项式），$\lim_{x \to -\infty} f(x) = -\infty$，$\lim_{x \to +\infty} f(x) = +\infty$；

由零点定理可知，方程$f(x) = 0$在$(-\infty, +\infty)$内\textbf{至少有一个实根}，排除 (A)。

\textbf{[2]} 求导与二次型换元：


\begin{equation*}
f'(x) = 5x^4 + 6ax^2 + 3b
\end{equation*}


令$t = x^2 \ge 0$，导函数化为关于$t$的一元二次函数$g(t) = 5t^2 + 6at + 3b$。

\textbf{[3]} 判别式判定导数恒正：

考察该二次方程的判别式$\Delta$：


\begin{equation*}
\Delta = (6a)^2 - 4(5)(3b) = 36a^2 - 60b = 12(3a^2 - 5b)
\end{equation*}


已知$3a^2 - 5b < 0$，故$\Delta < 0$恒成立。

因为二次项系数$5 > 0$且$\Delta < 0$，所以对任意实数$x$，恒有$f'(x) > 0$。

\textbf{[4]} 单调性结论：

导数恒大于0说明$f(x)$在$(-\infty, +\infty)$上严格单调增加，因此原方程\textbf{有且仅有唯一的实根}。

\textbf{[5]} 故正确选项为 \textbf{(B)}。
\end{solbox}


\subsection{例题 16: 试根法结合高阶导数不为零锁定根上限}

证明方程$2^x - x^2 - 1 = 0$有且仅有三个实根。

\textbf{\color{DeepPurple}【思路点拨】} 观察试根法找出3个确定的实根，再连续三次求导发现$f'''(x) > 0$恒成立，由罗尔定理推论锁死根的个数上限至多为3


操作[1]:存在性（找出3个确定的实根）：

\begin{solbox}[分步推导与答题规范]
记$f(x) = 2^x - x^2 - 1$：

- 代入$x = 0$：$f(0) = 2^0 - 0^2 - 1 = 0 \implies x_1 = 0$为第一个实根；
- 代入$x = 1$：$f(1) = 2^1 - 1^2 - 1 = 0 \implies x_2 = 1$为第二个实根；
- 考察区间$(2, 5)$：
 - $f(2) = 2^2 - 2^2 - 1 = -1 < 0$；
 - $f(5) = 2^5 - 5^2 - 1 = 32 - 25 - 1 = 6 > 0$。
 由零点定理，在$(2, 5)$内必存在$x_3$使得$f(x_3) = 0$。

由此确认方程至少存在3个不同实根：$x_1 = 0, x_2 = 1, x_3 \in (2, 5)$。
\end{solbox}

操作[2]:逐阶求导：

\begin{solbox}[分步推导与答题规范]
连续求导：

- $f'(x) = 2^x \ln 2 - 2x$
- $f''(x) = 2^x (\ln 2)^2 - 2$
- $f'''(x) = 2^x (\ln 2)^3$

由于底数$2 > 0$且$\ln 2 > 0$，故对任意实数$x$恒有$f'''(x) > 0$（即方程$f'''(x) = 0$无实根）。
\end{solbox}

操作[3]:由罗尔定理推论锁定根上限：

\begin{solbox}[分步推导与答题规范]
若方程$f'''(x) = 0$无实根（0个根），则原方程$f(x) = 0$在实数域内至多有$0 + 3 = 3$个实根。
\end{solbox}

操作[4]:结论：

\begin{solbox}[分步推导与答题规范]
综合可知，方程$2^x - x^2 - 1 = 0$\textbf{有且仅有三个实根}。
\end{solbox}


\subsection{例题 17: 参变量分离与动静态曲线交点分析}

已知方程$e^x = kx$有且仅有一个实根，则$k$的取值范围是\_\_\_\_\_\_。

\textbf{\color{DeepPurple}【思路点拨】} 排除$x=0$后分离参数$k = \frac{e^x}{x}$，转化为水平直线$y = k$与定曲线$g(x) = \frac{e^x}{x}$交点个数问题，由极值与渐近线求得范围


操作[1]:排除$x = 0$并分离参数：

\begin{solbox}[分步推导与答题规范]
若$x = 0$，左边$e^0 = 1 \neq 0$，故$x = 0$必非原方程的解。

两边同除以$x$分离参数$k$：


\begin{equation*}
k = \frac{e^x}{x}
\end{equation*}


原方程实根个数等价于水平直线$y = k$与定曲线$g(x) = \frac{e^x}{x}$的交点个数。
\end{solbox}

操作[2]:分析定曲线$g(x) = \frac{e^x}{x}$的性态：

\begin{solbox}[分步推导与答题规范]
定义域为$(-\infty, 0) \cup (0, +\infty)$。

- 一阶导数：$g'(x) = \frac{e^x(x - 1)}{x^2}$；
- 在$(0, +\infty)$上：驻点$x = 1$处取得极小值$g(1) = e$。当$x \in (0, 1)$时单减，当$x \in (1, +\infty)$时单增，正半轴值域为$[e, +\infty)$；
- 在$(-\infty, 0)$上：$g'(x) < 0$，$g(x) < 0$恒成立，值域为$(-\infty, 0)$。
\end{solbox}

操作[3]:水平直线$y = k$上下平移数交点：

\begin{solbox}[分步推导与答题规范]
- 当$k < 0$时：直线与负半轴单调分支\textbf{有且仅有1个交点}；
- 当$k = 0$时：无交点；
- 当$0 < k < e$时：无交点；
- 当$k = e$时：与正半轴曲线在极小值点$(1, e)$处相切，\textbf{有且仅有1个交点}；
- 当$k > e$时：在正半轴产生2个交点。
\end{solbox}

操作[4]:结论：

\begin{solbox}[分步推导与答题规范]
使得方程有且仅有一个实根的范围为 \textit{}$k = e$ 或 $k < 0$\textit{}。
\end{solbox}


\subsection{例题 18: 若尔当(Jordan)不等式证明(单调性法与凹凸性法)}

证明：当$0 < x < \frac{\pi}{2}$时，恒有$\sin x > \frac{2x}{\pi}$。

\textbf{\color{DeepPurple}【思路点拨】} 

方法一：同除以$x$转化为证明函数$f(x) = \frac{\sin x}{x}$在$\left(0, \frac{\pi}{2}\right]$上严格单调减少；
方法二：设差函数$F(x) = \sin x - \frac{2x}{\pi}$，利用二阶导数恒负说明曲线为严格凸弧，利用凸弧落在弦上方完成证明

操作[1]:方法一（单调性法）：

\begin{solbox}[分步推导与答题规范]
不等式两端同除以正数$x$，等价于证明$f(x) = \frac{\sin x}{x} > \frac{2}{\pi}$。

求导分析：


\begin{equation*}
f'(x) = \frac{x\cos x - \sin x}{x^2}
\end{equation*}


考察分子函数$h(x) = x\cos x - \sin x$：


\begin{equation*}
h'(x) = \cos x - x\sin x - \cos x = -x\sin x
\end{equation*}


当$x \in \left(0, \frac{\pi}{2}\right)$时，$x > 0$且$\sin x > 0$，故$h'(x) < 0$。

说明$h(x)$在$\left[0, \frac{\pi}{2}\right]$上严格单调减少，故$h(x) < h(0) = 0$。

代回得$f'(x) < 0$，故$f(x) = \frac{\sin x}{x}$在$\left(0, \frac{\pi}{2}\right]$上严格单调减少。

取右端点比较：


\begin{equation*}
f(x) > f\left(\frac{\pi}{2}\right) = \frac{\sin\frac{\pi}{2}}{\frac{\pi}{2}} = \frac{2}{\pi} \implies \sin x > \frac{2x}{\pi}
\end{equation*}

\end{solbox}

操作[2]:方法二（二阶导数凹凸性法）：

\begin{solbox}[分步推导与答题规范]
设辅助函数$F(x) = \sin x - \frac{2x}{\pi}\ \left(x \in \left[0, \frac{\pi}{2}\right]\right)$。

求二阶导数：


\begin{equation*}
F'(x) = \cos x - \frac{2}{\pi}, \quad F''(x) = -\sin x
\end{equation*}


当$x \in \left(0, \frac{\pi}{2}\right)$时，$F''(x) < 0$恒成立，曲线$y = F(x)$在区间上为\textbf{严格凸弧}。

计算端点函数值：


\begin{equation*}
F(0) = \sin 0 - 0 = 0, \quad F\left(\frac{\pi}{2}\right) = \sin\frac{\pi}{2} - \frac{2}{\pi} \cdot \frac{\pi}{2} = 1 - 1 = 0
\end{equation*}


由凸弧的几何特征（曲线上连接两个零点的弧段恒位于割线上方）：

在开区间$\left(0, \frac{\pi}{2}\right)$内部恒有$F(x) > 0$，即$\sin x > \frac{2x}{\pi}$。
\end{solbox}


\subsection{例题 19: 导函数极限趋零与严格单调性证明不等式}

证明：当$x > 0$时，恒有$\ln\frac{1+x}{x} - \frac{1}{1+x} < \frac{1}{\sqrt{x(1+x)}}$。

\textbf{\color{DeepPurple}【思路点拨】} 移项作差设函数$g(x)$，通过求导通分利用完全平方差公式判定$g'(x) > 0$，说明函数严格单增，再结合$x \to +\infty$时极限为零推导$g(x) < 0$


操作[1]:设差函数：

\begin{solbox}[分步推导与答题规范]
将欲证不等式移项，令：


\begin{equation*}
g(x) = \ln\left(1 + \frac{1}{x}\right) - \frac{1}{1+x} - \frac{1}{\sqrt{x(1+x)}} = \ln(1+x) - \ln x - \frac{1}{1+x} - [x(1+x)]^{-\frac{1}{2}}
\end{equation*}


只需证明当$x > 0$时，恒有$g(x) < 0$。
\end{solbox}

操作[2]:求导与通分：

\begin{solbox}[分步推导与答题规范]
对$g(x)$关于$x$求导：


\begin{equation*}
g'(x) = \frac{1}{1+x} - \frac{1}{x} + \frac{1}{(1+x)^2} - \left(-\frac{1}{2}\right)[x(1+x)]^{-\frac{3}{2}} (2x+1)
\end{equation*}



\begin{equation*}
= \frac{-1}{x(1+x)} + \frac{1}{(1+x)^2} + \frac{2x+1}{2[x(1+x)]^{\frac{3}{2}}} = \frac{-1}{x(1+x)^2} + \frac{2x+1}{2[x(1+x)]^{\frac{3}{2}}}
\end{equation*}


通分合并分子：


\begin{equation*}
g'(x) = \frac{-2\sqrt{x(1+x)} + (2x+1)}{2x(1+x)\sqrt{x(1+x)}} = \frac{2x + 1 - 2\sqrt{x^2+x}}{2[x(1+x)]^{\frac{3}{2}}}
\end{equation*}

\end{solbox}

操作[3]:判定导数符号：

\begin{solbox}[分步推导与答题规范]
考察分子各项平方差：


\begin{equation*}
(2x+1)^2 - \left(2\sqrt{x^2+x}\right)^2 = (4x^2 + 4x + 1) - (4x^2 + 4x) = 1 > 0
\end{equation*}


因为$2x+1 > 0$，所以$2x + 1 > 2\sqrt{x^2+x}$恒成立，分子恒大于0。

又分母在$x > 0$时显然大于0，因此对任意$x > 0$恒有$g'(x) > 0$。
\end{solbox}

操作[4]:求极限与得出结论：

\begin{solbox}[分步推导与答题规范]
由于$g'(x) > 0$，$g(x)$在$(0, +\infty)$上严格单调增加。

考察当$x \to +\infty$时的极限：


\begin{equation*}
\lim_{x \to +\infty} g(x) = \lim_{x \to +\infty} \left[\ln\left(1 + \frac{1}{x}\right) - \frac{1}{1+x} - \frac{1}{\sqrt{x(1+x)}}\right] = 0 - 0 - 0 = 0
\end{equation*}


一个在$(0, +\infty)$上严格单调增加且当$x \to +\infty$时极限趋近于0的函数，在其内部恒满足：


\begin{equation*}
g(x) < \lim_{x \to +\infty} g(x) = 0
\end{equation*}


由此证毕：$\ln\frac{1+x}{x} - \frac{1}{1+x} < \frac{1}{\sqrt{x(1+x)}}$。
\end{solbox}

\end{document}
